\section{Conclusiones}
\textbf{CONCLUSIÓN PRIMERA PARTE}
En la experiencia se observa una coherencia en los cálculos del factor de acoplamiento tomando el devanado primario como bobina de entrada como tomando el devanado secundario como bobina de entrada. Esto resulta coherente con el marco teórico pues el factor de acoplamiento es una propiedad intrínseca de la construcción del transformador y no debería cambiar en función de la bobina de entrada. Sim embargo, se observó un k mucho mas bajo del esperado. Como consecuencia, hay una mayor cantidad de flujo magnétco que circula por el núcleo de ferrite lo cual hace que el comportamiento del transformador analizado se aleje en mayor medida al comportamiento ideal. Lo demás resultados acerca de las inductancias de cada devanado resultan coherentes con las relaciones de tensiones entre ambos devanados. 

\subsection{Conclusi\'on del ensayo de aislamiento electrico}

A partir de las mediciones realizadas, se verifico el efecto de aislamiento
electrico del transformador. Al variar el offset de entrada entre
\(-5\,\mathrm{V}\), \(0\,\mathrm{V}\) y \(+5\,\mathrm{V}\), la tension medida del
lado de entrada modifico su valor maximo, pero la tension de salida del
transformador no presento variaciones apreciables.

Esto confirma que no existe una conexi\'on el\'ectrica directa entre ambos
bobinados y que la transferencia de energ\'ia se realiza mediante acoplamiento
magn\'etico. En consecuencia, la componente continua de la se\~nal no se
transmite al secundario, mientras que la componente alterna s\'i aparece en la
salida. Las peque\~nas diferencias observadas pueden atribuirse a la resoluci\'on
de los instrumentos, a la tolerancia de \(R_m\), a la impedancia interna de la
fuente y a errores propios de lectura en el osciloscopio.

\textbf{CONCLUSIÓN SEGUNDA PARTE}
Los resultados del cálculo de las impedancias de entrada del transformador con carga resultan coherentes con las simulaciones. La impedancia que difiera más con el resultado simulado es la correspondiente a tomar el segundo devanado como bobina de entrada. Este error es esperable pues como es el lado de alta tensión, la impedancia vista por la fuente resulta muy grande dando como resultado una corriente muy pequeña. La medición de esta corriente está sujeta a errores por parte del instrumento de medición. 